La fonction nommée $C$ exprime la concentration en quantité de matière en composé P en fonction de $t$, exprimé en secondes, dans l'intervalle $[0 ; 1200]$. Elle est définie par :
\[C(t)=3,1 \times 10^{-5} \times \e^{-0,00702 t}\]
La fonction $C$ est dérivable sur l'intervalle $[0 ; 1200]$.
\begin{enumerate}[start=7]
	\item Déterminer la fonction dérivée de $C$ notée $C'$ ou $\dfrac{\mathrm{d} C}{\mathrm{~d} t}$ sur l'intervalle $[0; 1200]$.
	\stepcounter{enumi}

	\item Résoudre sur $[0 ; 1200]$ l'équation
	\[3,1 \times 10^{-5} \times \e^{-0,00702 t}=3,1 \times 10^{-6}\]
	Arrondir le résultat à l'unité.
\end{enumerate}


\needspace{10\baselineskip}
% saut de page si moins de 10 lignes disponibles
